% ===========================================================================
\section{Two pillars: $\KL\ge0$ and the compression limit}\label{sec:source}
\warmth{0}\quad\whisper{One inequality implies all the bounds; one theorem turns entropy into bits on a disk.}

\begin{strand}
The whole edifice rests on two results. The first is an \emph{inequality} from which every other
bound drops out: relative entropy is never negative. The second is the \emph{operational} payoff
that makes the bits real: Shannon's source coding theorem, which says the average code length can
get down to the entropy, and no further.
\end{strand}

\block{Pillar 1: nothing beats the truth}

\begin{theorem}[Gibbs' inequality]\label{thm:gibbs}
$\KL(p\,\Vert\,q)\ge 0$ for all distributions $p,q$, with equality iff $p=q$.
\end{theorem}
\begin{proof}
Apply Jensen to the concave $\log$:
$-\KL(p\Vert q)=\sum_x p(x)\log\frac{q(x)}{p(x)}\le\log\sum_x p(x)\frac{q(x)}{p(x)}=\log 1=0$.
\TODO{state the equality condition from strict concavity, and derive the two corollaries below}
\end{proof}

\recall{Two corollaries, for free, by choosing $q$: take $q$ uniform to get $\Ent(X)\le\log\abs{\X}$ (uniform is maximally uncertain); take $q=p_Xp_Y$ to get $\MI(X;Y)\ge0$, i.e.\ $\Ent(X\mid Y)\le\Ent(X)$.}

\block{Pillar 2: entropy is the floor on code length}

\begin{theorem}[Kraft + source coding]\label{thm:source}
Any binary \keyword{prefix} code with lengths $\ell(x)$ obeys $\sum_x 2^{-\ell(x)}\le 1$ (Kraft).
Consequently the expected length $L=\sum_x p(x)\ell(x)$ satisfies
\[
  L\;\ge\;\Ent(X),
\]
and a code achieving $L<\Ent(X)+1$ always exists (take $\ell(x)=\lceil-\log p(x)\rceil$).
\end{theorem}
\begin{proof}
For the lower bound, $L-\Ent(X)=\sum_x p(x)\log\frac{p(x)}{2^{-\ell(x)}}=\KL(p\Vert r)+\log\frac1{\sum 2^{-\ell}}\ge0$,
where $r(x)=2^{-\ell(x)}/\sum 2^{-\ell}$. \TODO{check both terms are $\ge0$ using Kraft and Gibbs}
\end{proof}

\begin{yourturn}
Code the distribution $p=\big(\tfrac12,\tfrac14,\tfrac18,\tfrac18\big)$.
\[
  \Ent(X)=\fillin[1.6cm]\text{ bits},\qquad L=\fillin[1.6cm]\text{ bits}.
\]
\[
  \text{optimal code lengths}\ (\ell_1,\ell_2,\ell_3,\ell_4)=\fillin[2.6cm].
\]
\recall{$\Ent=1.75$ bits; optimal lengths $(1,2,3,3)$; $L=\Ent$ exactly because the probabilities are \keyword{dyadic} (powers of two). Otherwise you pay up to one extra bit.}
\end{yourturn}

\begin{strand}
Squeeze out the last bit by coding \emph{blocks}: over $n$ symbols the per-symbol overhead is
$1/n\to0$, so the rate $\to\Ent(X)$. The clean statement is the AEP / typical set: about
$2^{n\Ent}$ sequences carry essentially all the probability, so $n\Ent$ bits suffice (§\ref{sec:faces}).
\end{strand}
